% -------------------------------------------------------------------------
% WBS ID: RES-OPT-CST
% Deliverable: Constraints, Feasibility and Design Admissibility
% Status: In progress
% Evidence status: Optimisation constraints deferred; geometry feasibility linked
% Review status: Not reviewed
% Last updated: 2026-07-19
% Dependencies: RES-GEO-CST, RES-VER-SPEC
% Downstream interfaces: Future RES-OPT-SPEC
% -------------------------------------------------------------------------

\section{RES-OPT-CST --- Constraints, Feasibility and Design Admissibility}
\label{sec:res-opt-cst}

\subsection{Purpose and Scope}
\label{sec:res-opt-cst-purpose}

This Deliverable records future optimisation constraints as deferred.
The active project uses RES-GEO-CST feasibility checks for fixed-rigid
barrier-class comparison; it does not define an admissible
optimisation space or reinforcement-intervention campaign.

\subsection{Research Questions}
\label{sec:res-opt-cst-research-questions}

% TODO: State the questions that must be answered before the Deliverable
% can be accepted.

\subsection{Terminology and Definitions}
\label{sec:res-opt-cst-definitions}

% TODO: Define the technical terminology, symbols, quantities and
% classifications used in this Deliverable.

\subsection{Literature Evidence}
\label{sec:res-opt-cst-literature-evidence}

% TODO: Record evidence from peer-reviewed papers, authoritative datasets,
% standards, technical reports and primary documentation.
%
% Do not create a detached literature-summary list. Explain how each source
% supports, contradicts or limits a project decision.

\subsection{Governing Relations and Quantitative Definitions}
\label{sec:res-opt-cst-governing-relations}

% TODO: Record relevant equations, dimensionless groups, empirical models,
% extraction rules or quantitative definitions.
%
% Use align environments for displayed equations.
% Every displayed equation must have a unique \label{eq:...}.
% Do not place punctuation after displayed equations.

\subsection{Decision Variables}
\label{sec:res-opt-cst-decision-variables}

% TODO: Complete this RES-OPT Domain-specific subsection.

\subsection{Objective Functions}
\label{sec:res-opt-cst-objective-functions}

% TODO: Complete this RES-OPT Domain-specific subsection.

\subsection{Constraints}
\label{sec:res-opt-cst-constraints}

Future constraints may include geometry validity, constructability,
material limits, foundation limits, meshability, numerical feasibility,
validated-domain limits and direct-simulation confirmation
requirements. In the active scope, only geometry validity,
terrain/corridor fit, meshability and unresolved-physics flags are
required for barrier-class comparison.

\subsection{Sampling Strategy}
\label{sec:res-opt-cst-sampling-strategy}

% TODO: Complete this RES-OPT Domain-specific subsection.

\subsection{Optimisation Method}
\label{sec:res-opt-cst-optimisation-method}

% TODO: Complete this RES-OPT Domain-specific subsection.

\subsection{Computational Budget}
\label{sec:res-opt-cst-computational-budget}

% TODO: Complete this RES-OPT Domain-specific subsection.

\subsection{Termination and Selection Criteria}
\label{sec:res-opt-cst-termination-selection-criteria}

% TODO: Complete this RES-OPT Domain-specific subsection.

\subsection{Candidate Approaches}
\label{sec:res-opt-cst-candidate-approaches}

% TODO: Define the credible candidate approaches considered.

\subsection{Critical Comparison}
\label{sec:res-opt-cst-critical-comparison}

% TODO: Compare candidates using physically and project-relevant criteria.
% Do not select an approach solely on popularity or implementation ease.

\subsection{Assumptions and Applicability}
\label{sec:res-opt-cst-assumptions}

% TODO: State assumptions and the conditions under which they are valid.

\subsection{Limitations and Failure Conditions}
\label{sec:res-opt-cst-limitations}

% TODO: State where the evidence, model or selected approach ceases to be
% reliable or applicable.

\subsection{Project-Specific Conclusions}
\label{sec:res-opt-cst-conclusions}

% TODO: Convert the evidence into conclusions specific to the Studentship
% project. Do not merely restate literature findings.

\subsection{Selected Approach and Rationale}
\label{sec:res-opt-cst-selected-approach}

% TODO: Record the selected approach only after sufficient evidence exists.
% State the alternatives rejected and the reasons for rejection.

\subsection{Unresolved Decisions}
\label{sec:res-opt-cst-unresolved-decisions}

% TODO: Record unresolved questions, required evidence and decision owners.

\subsection{Requirements and Handoffs}
\label{sec:res-opt-cst-requirements}

% TODO: State requirements passed to other Research Domains, Software,
% verification activities or later execution work.

\subsection{Deliverable Completion Record}
\label{sec:res-opt-cst-completion-record}

% TODO: Record whether the Deliverable is:
%
% - Not started
% - In progress
% - Draft complete
% - Reviewed
% - Accepted
%
% Acceptance requires:
%
% - the research questions to be answered;
% - claims to be supported by appropriate evidence;
% - assumptions and limitations to be explicit;
% - the project-specific conclusion to be stated;
% - downstream requirements to be traceable;
% - unresolved decisions to be closed or formally deferred.
